\section{Source locators and conventions for Table~\ref{tab:resource-comparison}}
\label{app:oracle-comparison}

This appendix is a verification aid for Table~\ref{tab:resource-comparison},
not a separate survey.  For each external row, it records the linked version,
the theorem or definition supporting the entry, and only the conversion needed
to read the table.  Numbering below refers to the linked version and may differ
from another bibliographic version.

\paragraph{Common conventions.}
For an external randomized schedule $M$ with conditional output law $Q_M$,
\emph{Mean KL/TV} denotes $\mathbb E_M D(P,Q_M)$, not
$D(P,\mathbb E_MQ_M)$; an unqualified TV guarantee concerns the randomized
output law.  The marker \emph{exp.} applies to both resource
columns, while \emph{same} means that the submission bound also bounds depth.
Appendix~\ref{app:comparison-states} gives the conversion from conditional-row
queries or unmasking updates to our masked-state submissions.  Fixed accuracy
and polynomial vocabulary are imposed only when comparing growth in $N$.
Uniform counterfactual-row error and an average under a data--mask or trajectory
law are different oracle assumptions; Table~\ref{tab:resource-comparison}
does not identify them.

\subsection{Distribution-general sampler rows}
\label{app:comparison-exact-oracles}

\begin{description}
\item[\Citet{anari2024parallel}.]
\href{https://arxiv.org/pdf/2408.09442v1}{arXiv:2408.09442v1},
Theorem~2 and Algorithm~4: exact conditional rows and exact output, with
$O(N)$ expected row queries and $\widetilde O(N^{2/3})$ expected depth for
polynomial vocabulary.  Appendix~\ref{app:comparison-states} supplies the
submission conversion used in the table.

\item[\Citet{anari2026autospeculation}.]
\href{https://arxiv.org/pdf/2511.07869v1}{arXiv:2511.07869v1},
Theorem~27 and footnote~3: $O(N\log N)$ expected full-row queries and
$O(\sqrt{N\log\VocabSize\,(\log N)^3})$ expected depth.
Remark~32, equation~(25), gives the noisy-row TV entry:
\begin{equation}
 e\le\min\{\tau/N,N^{-3/2}\}
 \quad\Longrightarrow\quad
 \mathrm{TV}(P,Q)\le\tau.
 \label{app:eq:autospec-native-tv}
\end{equation}
The condition is uniform TV accuracy of the normalized rows; the conclusion is
TV for the output law.  Exact rows give exact output.
\end{description}

\subsection{Dependence-based schedule rows}
\label{app:comparison-unmasking-errors}

\begin{description}
\item[\Citet{li2025convergence}.]
\href{https://arxiv.org/html/2505.21400v2}{arXiv:2505.21400v2},
Definition~1, equation~(7), defines the predictor error averaged over target
data, random masks, and the weighted time index.  Theorem~1, equation~(8), and
Corollary~1, equation~(9), give for near-balanced blocks
\[
 \mathbb E_M\mathrm{KL}(P\Vert Q_M)
 \le \frac{C}{K}\bigl(\mathrm{TC}(P)+\mathrm{DTC}(P)\bigr)
      +\varepsilon_{\rm train},
\]
where $M$ is the random schedule and $K$ is its number of updates.  The table
allocates a constant accuracy budget to $\varepsilon_{\rm train}$ and
substitutes the path calculation in
Proposition~\ref{app:prop:path-in-class}; no uniform counterfactual-row
condition is inferred from this average.

\item[\Citet{chen2026optimal}.]
\href{https://arxiv.org/html/2511.04647v2}{arXiv:2511.04647v2},
Definition~1.1 specifies schedule-averaged forward KL, Definition~2.1 defines
one partial-assignment query returning all remaining conditional rows, and
Theorem~1.9 gives, for a supplied TC or DTC upper bound $\widehat T$,
\begin{equation}
 \mathbb E_M\mathrm{KL}(P\Vert Q_M)\le\varepsilon,
 \qquad
 K\le2+(1+\log N)(1+\lceil\widehat T/\varepsilon\rceil).
 \label{app:eq:chen-native-bound}
\end{equation}
The table uses exact rows and a constant-factor supplied bound.  On the path of
Proposition~\ref{app:prop:path-in-class}, both TC and DTC are $\Theta(N)$.
\end{description}

\paragraph{Baselines and this paper.}
Exact singleton sampling uses $N$ submissions and depth $N$ by the chain rule.
For marker $a$, one exact product batch has forward KL
$\mathrm{TC}(P)$, which is linear on the path of
Proposition~\ref{app:prop:path-in-class}; it is only a resource baseline.
Our two rows are the balanced bounds in
Theorem~\ref{main:thm:upper-curve} and
\eqref{main:eq:balanced-upper}: case (i) uses uniform row squared-Hellinger
error and reports mean TV, while case (ii) uses uniform row forward KL and
reports mean forward KL.  Their exact finite statements are
Theorems~\ref{app:thm:finite-upper} and
\ref{app:thm:finite-kl-upper}.

\subsection{Supplementary citation map}
\label{app:related-work-supplement}

The following citations supply context only; they are not used to derive a row
of Table~\ref{tab:resource-comparison}.

{\color{red}
\paragraph{Modeling background.}
\begin{itemize}
\item \emph{Forest approximation.}
\citet{liu2011forestdensity} study nonparametric density estimation using
forests, allowing the true continuous distribution to lie outside that
family (\href{https://jmlr.org/papers/v12/liu11a.html}{abstract}).
This motivates forests as an approximation class; our target instead
satisfies exact forest factorization.

\item \emph{Prediction without exact structure recovery.}
\citet{bresler2020predictions} study tree Ising prediction using
\emph{small-set TV}, which compares marginals on small subsets
(\href{https://arxiv.org/abs/1604.06749}{abstract}).
This differs from the global TV objective of
\citet{daskalakis2021treeising} and our case (i).

\item \emph{Learned conditional distributions.}
\citet{heckerman2000dependency} allow separately learned local conditionals
that need not be compatible with a joint law
(\href{https://www.jmlr.org/papers/volume1/heckerman00a/heckerman00a.pdf\#page=7}{Section 3,
pp.~55--56}). Their inference procedure uses ordered pseudo-Gibbs updates;
our interface supplies masked-state rows for irreversible commits.
\end{itemize}

\paragraph{Algorithmic antecedents.}
\label{app:algorithmic-antecedents}
The construction in Section~\ref{sec:main-algorithm} connects the following
design principles to hidden-forest sampling.
\begin{itemize}
\item \emph{Shared evaluations (Steps 2--3).}
\citet{coleman1983estimation} use graph coloring to reduce function
evaluations for sparse Jacobian estimation (see the paper's
\href{https://epubs.siam.org/doi/10.1137/0720013}{abstract}).
Our packed probes share evaluations across source--readout pairs;
the separation conditions (R)/(S) in Section~\ref{sec:main-packed-screen}
make their conditional replies reproduce single-source tests.

\item \emph{Hashing and majority recovery (Step 3).}
\citet{bshouty2018exact},
\href{https://arxiv.org/pdf/1706.06934v1\#page=17}{arXiv v1, Section 4.3,
Theorem 8, Figure 1, steps 1--3}, hash variables, lift identified buckets
back to variables, and aggregate by majority.
Our bucket tests use conditional-row diameters; we recover low-degree
neighborhoods rather than a Boolean function.

\item \emph{Peeling (Step 4).}
\citet{becker2010referee},
\href{https://arxiv.org/pdf/1009.4447v2\#page=9}{arXiv v2, Section 3.1, p.~9},
reconstruct a forest by decoding leaves from exact degrees and neighbor-ID
sums, then removing them and updating these summaries.
Here, low-degree reports instead select singleton commits that shrink
the high-degree core (\cref{main:eq:peeling-contraction}), without requiring
complete neighborhoods for the committed vertices.

\item \emph{Separator layers (Steps 5--6).}
\citet{iyer1988optimal} relate vertex ranking to separator-tree height
(see their
\href{https://research.ibm.com/publications/optimal-node-ranking-of-trees}{abstract}).
In our centroid schedule, earlier layers have higher ranks: their committed
vertices separate same-layer vertices into distinct residual components.
We use centroid recursion, not their optimal-ranking algorithm.
\end{itemize}
}

\paragraph{\textcolor{red}{Additional context.}}
Diffusion foundations and discrete or language formulations include
\citet{sohldickstein2015deep,ho2020denoising,song2021score,
campbell2022continuous,li2022diffusionlm,he2023diffusionbert}.
Parallel-time and complementary complexity perspectives include
\citet{shih2023parallel,chen2024sublinear,yao2026parallelintime,
feng2025theoretical,jiang2026optimal,zhang2026generation,liang2026sharp}.

Adaptive order, lookahead, confidence, and schedule analyses include
\citet{li2024promises,kim2025train,hayakawa2026demystifying,
lee2025lookahead,fu2025bits,lavenant2025error,zhao2026adaptation,
cai2026confidence,jys}.
For the specific round statement cited in the main text, the locator for
\citet{fu2025bits} is
\href{https://arxiv.org/html/2511.21103v1}{arXiv:2511.21103v1},
Assumption~3.1 and Theorem~3.2.
Dependence-based, remasking, and within-batch modeling perspectives include
\citet{dmitriev2026efficient,dmitriev2026remasking,wainwright2026geometry,
bansal2026randomwalks,hayakawa2025distillation,liu2024discrete,
lezama2022discrete}.

Related oracle, learning, and convergence settings include approximate-density
and conditional-sampling access
\citep{golowich2026testtime,canonne2015conditional,canonne2021random,
chen2021junta,blanca2023coordinate}, distributed sampling
\citep{feng2020local}, learned-denoiser error
\citep{wakasugi2025state}, and discrete-diffusion convergence
\citep{zhang2025convergence,ren2025discrete,liang2025absorb}.
These works use objectives or access models different from the fixed
singleton-conditional interface studied here; no quantitative conversion is
claimed beyond the four source rows above.

\subsection{Reading Table~\ref{tab:resource-comparison}}
\label{app:comparison-table-notes}

\begin{itemize}
\item The external entries are sufficient upper bounds, not lower bounds on
their samplers.  The two Anari rows retain their distribution-general
guarantees.  For the TC/DTC rows, Appendix~\ref{app:comparison-path}
substitutes one path with
$\mathrm{TC}(P)=\mathrm{DTC}(P)=\Theta(N)$; this does not imply that those
samplers require linear work.

\item Marker $b$ is exactly the pair of row-precision requirements in
\eqref{app:eq:autospec-native-tv}.  For
\citet{chen2026optimal}, $\widehat T$ is a supplied TC or DTC upper bound.
For the Li--Cai row, $K=N$ removes the factorization term in its Theorem~1;
for the Chen et al. row, exact singleton sampling gives the $N$-call cap.

\item The scaling comparison fixes positive output tolerance and polynomial
vocabulary.  Constant allocations to source-specific learning or sampling
error do not change the displayed powers of $N$.

\item Our rows use $\RankTailExponent>1$ and the balanced examples of
\eqref{main:eq:balanced-upper}.  Their common leading resource power is
$2/3+1/(9\RankTailExponent)<1$; the target class and oracle assumptions remain
part of the comparison.
\end{itemize}
